% Part: many-valued-logic
% Chapter: syntax-and-semantics
% Section: semantic-notions

\documentclass[../../../include/open-logic-section]{subfiles}

\begin{document}

\olfileid{mvl}{syn}{sem}

\olsection{অর্থগত ধারণাসমূহ}

ধরা যাক, বহুমানী যুক্তিবিদ্যা~$\Log L$ একটি ম্যাট্রিক্স দিয়ে দেওয়া আছে।
তাহলে আমরা~$\Log L$-এর জন্য প্রচলিত অর্থগত ধারণাগুলি সংজ্ঞায়িত করতে পারি।

\begin{defn}
\begin{enumerate}
\item কোনো !!^a{formula}~$!A$ \emph{পরিতৃপ্তিযোগ্য}, যদি কোনো
  $\pAssign{v}$-এর জন্য $\pSat{v}{!A}$ হয়; কোনো $\pAssign{v}$-এর
  জন্যই $\pSat{v}{!A}$ না হলে সেটি \emph{অপরিতৃপ্তিযোগ্য};
\item কোনো !!^a{formula}~$!A$ \emph{সর্বতঃসত্য}, যদি সব
  !!{valuation}s~$v$-এর জন্য $\pSat{v}{!A}$ হয়;
\item $\Gamma$ যদি !!{formula}s-এর একটি সেট হয়, তবে
  $\Gamma \Entails !A$ (“$\Gamma$ থেকে $!A$ অনুসৃত হয়”) তখনই, যখন
  $\pSat{v}{\Gamma}$ হয় এমন প্রত্যেক !!{valuation}~$\pAssign{v}$-এর
  জন্য $\pSat{v}{!A}$ হয়।
\item $\Gamma$ যদি !!{formula}s-এর একটি সেট হয়, তবে এমন
  !!a{valuation}~$\pAssign{v}$ থাকলে $\Gamma$ \emph{পরিতৃপ্তিযোগ্য},
  যার জন্য $\pSat{v}{\Gamma}$; অন্যথায় $\Gamma$
  \emph{অপরিতৃপ্তিযোগ্য}।
\end{enumerate}
\end{defn}

ধ্রুপদি যুক্তিবিদ্যার ক্ষেত্রে এই ধারণাগুলির জন্য যে তথ্যগুলি পাওয়া যায়,
সেগুলির কয়েকটি এখানেও পাওয়া যায়:

\begin{prop}
\ollabel{prop:semanticalfacts}
\begin{enumerate}
\item $!A$ তখনই সর্বতঃসত্য, যখন
  $\emptyset \Entails !A$;
\item $\Gamma$ পরিতৃপ্তিযোগ্য হলে $\Gamma$-র প্রতিটি সসীম উপসেটও
  পরিতৃপ্তিযোগ্য;
\item\ollabel{def:monotonicity}%
একঘেয়েতা: যদি $\Gamma \subseteq \Delta$
  এবং $\Gamma \Entails !A$ হয়, তবে $\Delta \Entails !A$-ও হয়;
\item\ollabel{def:Cut}%
পরিযায়িতা: যদি $\Gamma \Entails !A$ এবং
  $\Delta \cup \{ !A\} \Entails !B$ হয়, তবে $\Gamma \cup \Delta \Entails
  !B$;
\end{enumerate}
\end{prop}

\begin{proof}
অনুশীলনী।
\end{proof}

\begin{prob}
\olref[mvl][syn][sem]{prop:semanticalfacts} প্রমাণ করো।
\end{prob}

ধ্রুপদি যুক্তিবিদ্যায় অনুসিদ্ধান্তের সঙ্গে শর্তবচনের সম্পর্ক স্থাপন করা
যায়। যেমন, \emph{মোডাস পোনেন্স}-এর বৈধতা আছে: যদি $\Gamma
\Entails !A$ এবং $\Gamma \Entails !A \lif !B$ হয়, তবে $\Gamma \Entails
!B$। ধ্রুপদি যুক্তিবিদ্যায় $\Entails$ ও~$\lif$-এর আরেকটি গুরুত্বপূর্ণ
সম্পর্ক হলো অর্থগত নিঃসরণ উপপাদ্য: $\Gamma \Entails !A
\lif !B$ তখনই, যখন $\Gamma \cup \{!A\} \Entails !B$। বহুমানী
যুক্তিবিদ্যায় এই ফলগুলি \emph{সবসময় সত্য নয়}। সেগুলি সত্য কি না,
তা সত্যমান-অপেক্ষক~$\tf{\lif}$-এর উপর নির্ভর করে।

\end{document}
